%==============================================================================
% CHAPTER 19: Chromatin Breathing and Splicing as Partition Selection
%==============================================================================

\chapter{Chromatin Breathing and Splicing as Partition Selection}
\label{chap:chromatin_splicing}

\epigraph{%
  What the ear hears not, the breath must carry.\\
  What the breath carries not, the structure must decide.%
}{after Wilhelm von Humboldt, \textit{On the Structural Variety of Human Language} (1836)}

\begin{chapterbox}
Nucleosome breathing --- the reversible partial unwrapping of DNA from the
histone octamer --- is driven by Debye length oscillations rather than by
ATP-dependent remodellers acting in isolation. The Charge-Driven Nucleosome
Breathing Theorem establishes that metabolic oscillations modulate nucleosome
stability through electrostatic potential, producing the observed 50\%
amplitude oscillation at metabolic timescales. A nucleosome is formally
identified as a partition gate: when closed it sets the partition depth of
the underlying DNA sequence to $\Depth_{\mathrm{access}} \to \infty$; when
open it reduces $\Depth_{\mathrm{access}}$ to allow transcription factor
binding. Alternative splicing is recast as categorical partition selection:
each exon inclusion or exclusion is a partition boundary crossing event, with
inclusion probability given by the Boltzmann-like formula
$P_i = b^{-\Depth_{\mathrm{include},i}}/\sum_j b^{-\Depth_{\mathrm{include},j}}$.
The metabolic coupling to splicing through SR protein phosphorylation explains
genome-wide splicing changes within minutes of metabolic perturbation.
Histone modifications and DNA methylation are identified as partition depth
increments --- the histone code is the partition code.
\end{chapterbox}

%------------------------------------------------------------------------------
\section{Nucleosome Breathing: Observations and Mechanisms}
\label{sec:breathing_overview}
%------------------------------------------------------------------------------

The nucleosome is the fundamental unit of eukaryotic chromatin. Approximately
147 base pairs of DNA wrap $1.65$ times around a histone octamer (two copies
each of H2A, H2B, H3, and H4), burying the underlying sequence from
solution. This wrapping is not permanent: single-molecule and FRET studies
demonstrate that nucleosome ends spontaneously unwrap and rewrap on timescales
of milliseconds to seconds --- a motion termed \emph{nucleosome breathing}.

Standard biochemical accounts attribute breathing to: (a) thermal fluctuations
in the histone--DNA contacts at the DNA entry/exit points; (b) ATP-dependent
remodellers (SWI/SNF, ISWI, CHD, INO80 families) that actively slide or
evict nucleosomes. Partition theory adds a third mechanism that is not
downstream of these two but upstream and more general: direct electrostatic
modulation by the metabolic Debye length oscillation.

%------------------------------------------------------------------------------
\section{The Electrostatic Balance of the Nucleosome}
\label{sec:nucleosome_electrostatics}
%------------------------------------------------------------------------------

The stability of the nucleosome is an electrostatic problem. DNA carries
$-2e$ per base pair; the 147\,bp contact patch carries
$-2 \times 147 = -294\,e$ in total. The histone octamer presents approximately
$+200$ positively charged amino acid residues (lysines and arginines) within
the DNA-binding surface. The net charge on the nucleosome core particle is
therefore approximately $-294 + 200 = -94\,e$ under physiological conditions,
with the octamer neutralising $\sim 68\%$ of the DNA charge.

The electrostatic binding energy between the wrapped DNA and the octamer is

\begin{equation}
  U_{\mathrm{bind}} = \frac{Q_{\mathrm{DNA}}\,Q_{\mathrm{histones}}}{4\pi\varepsilon_0\varepsilon_r\,r_{\mathrm{eff}}}
  \exp\!\left(-\frac{r_{\mathrm{eff}}}{\lambda_D}\right),
  \label{eq:nucleosome_binding_energy}
\end{equation}
where $Q_{\mathrm{DNA}} = -294\,e$, $Q_{\mathrm{histones}} = +200\,e$,
$r_{\mathrm{eff}} \approx 2$\,nm is the effective charge--charge separation
averaged over the DNA-wrap geometry, $\varepsilon_r \approx 80$ is the
effective dielectric constant, and $\lambda_D$ is the Debye length from
Eq.~\eqref{eq:debye_length} in Chapter~\ref{chap:metabolic_clock}.

\begin{theorem}[Charge-Driven Nucleosome Breathing]
\label{thm:nucleosome_breathing}
Metabolic oscillations drive nucleosome breathing with amplitude proportional
to the fractional oscillation in $\lambda_D$. Specifically:
\begin{enumerate}[label=(\roman*)]
  \item When $\lambda_D$ increases (lower ionic strength, less Debye
        screening): $|U_{\mathrm{bind}}|$ increases, the nucleosome is more
        stable, and chromatin is compact.
  \item When $\lambda_D$ decreases (higher ionic strength, more screening):
        $|U_{\mathrm{bind}}|$ decreases, the electrostatic attraction is
        weakened, the nucleosome breathes open, and chromatin is accessible.
\end{enumerate}
The fractional amplitude of nucleosome breathing oscillations is
\begin{equation}
  \frac{\delta\theta}{\bar\theta} = \frac{r_{\mathrm{eff}}}{\bar\lambda_D}\,
  \frac{\delta\lambda_D}{\bar\lambda_D} \cdot
  \frac{|U_{\mathrm{bind}}|}{\kB T},
  \label{eq:breathing_amplitude}
\end{equation}
where $\theta$ is the unwrapping angle (fraction of DNA length unwrapped).
\end{theorem}

\begin{proof}
Differentiate $U_{\mathrm{bind}}(\lambda_D)$ in Eq.~\eqref{eq:nucleosome_binding_energy}
with respect to $\lambda_D$:
\begin{equation}
  \frac{\partial U_{\mathrm{bind}}}{\partial\lambda_D} = U_{\mathrm{bind}} \cdot
  \frac{r_{\mathrm{eff}}}{\lambda_D^2}.
\end{equation}
A fractional change $\delta\lambda_D/\bar\lambda_D$ in Debye length induces
\begin{equation}
  \delta U_{\mathrm{bind}} = U_{\mathrm{bind}} \cdot \frac{r_{\mathrm{eff}}}{\bar\lambda_D}
  \cdot \frac{\delta\lambda_D}{\bar\lambda_D}.
\end{equation}
The equilibrium unwrapping angle $\theta$ satisfies the Boltzmann balance
$\theta/(1-\theta) = \exp(-U_{\mathrm{bind}}/\kB T)$. Linearising around
$\bar\theta$ gives $\delta\theta = \bar\theta(1-\bar\theta)\,\delta U_{\mathrm{bind}}/\kB T$.
For $\bar\theta \approx 0.5$ (50\% mean breathing amplitude as observed),
$\bar\theta(1-\bar\theta) = 0.25$, and the stated result follows. $\blacksquare$
\end{proof}

\begin{keyresult}
The observed 50\% amplitude of nucleosome breathing oscillations at metabolic
timescales is the quantitative signature of electrostatic coupling between the
Debye length oscillation and the histone--DNA binding energy. No additional
mechanism is required.
\end{keyresult}

%------------------------------------------------------------------------------
\section{The Nucleosome as a Partition Gate}
\label{sec:nucleosome_gate}
%------------------------------------------------------------------------------

\begin{theorem}[Nucleosome as Partition Gate]
\label{thm:nucleosome_gate}
A nucleosome positioned over a DNA sequence $\sigma$ implements a partition
gate on $\sigma$:
\begin{enumerate}[label=(\roman*)]
  \item \textbf{Closed gate (compact nucleosome):} The partition depth of
        $\sigma$ as seen from the nuclear solvent is
        $\Depth_{\mathrm{access}}(\sigma) \to \infty$. Transcription factors
        with finite affinity cannot access $\sigma$.
  \item \textbf{Open gate (breathing nucleosome):} The DNA ends of $\sigma$
        are exposed. The effective partition depth decreases to
        $\Depth_{\mathrm{access}}(\sigma) = \Depth_{\mathrm{bare}} +
        \Depth_{\mathrm{breathing}}$, where $\Depth_{\mathrm{bare}}$ is the
        intrinsic sequence-level partition depth and $\Depth_{\mathrm{breathing}}$
        is the residual barrier from partial wrapping.
\end{enumerate}
Chromatin remodellers (SWI/SNF, ISWI, CHD, INO80) are \emph{partition topology
modifiers}: they expend ATP to change the $\Depth$ landscape of a chromatin
region, either by sliding the gate (repositioning) or removing it (eviction).
\end{theorem}

\begin{proof}
The partition depth $\Depth$ of a sequence is defined as
$\Depth = \log_b(\Omega_{\mathrm{total}}/\Omega_{\mathrm{accessible}})$,
the logarithm of the ratio of total conformational states to accessible ones.
When a nucleosome is fully wrapped, the underlying DNA is sterically excluded
from solution; the number of conformational states accessible to a diffusing
transcription factor at the binding site approaches zero, so
$\Omega_{\mathrm{accessible}} \to 0$ and $\Depth \to \infty$. When the
nucleosome breathes open with unwrapping angle $\theta$, the accessible
states scale as $\theta\,\Omega_{\mathrm{open}}$, giving
$\Depth = \log_b(\Omega_{\mathrm{total}}/(\theta\,\Omega_{\mathrm{open}}))
= \Depth_{\mathrm{bare}} - \log_b\theta$.
Remodellers shift $\theta$ by hydrolyzing ATP to apply mechanical force to
the histone--DNA interface; they are therefore topological modifiers of the
$\Depth$ landscape. $\blacksquare$
\end{proof}

%------------------------------------------------------------------------------
\section{Alternative Splicing as Categorical Partition Selection}
\label{sec:splicing_partitions}
%------------------------------------------------------------------------------

Pre-mRNA splicing is the process by which introns are removed and exons are
joined to form mature mRNA. When a pre-mRNA contains multiple alternative
exons, the spliceosome must \emph{choose} which exons to include and which to
skip. This choice determines the protein isoform produced, and thus the
functional output of gene expression. The standard description frames the
choice as a competition between splice sites for spliceosomal recognition. The
partition description is more fundamental and more precise.

\begin{definition}[Partition Boundary in Pre-mRNA]
\label{def:splice_boundary}
Each splice site in a pre-mRNA is a partition boundary. Including exon $i$
in the mature mRNA corresponds to \emph{not crossing} the partition boundary
at the flanking splice sites; skipping exon $i$ corresponds to crossing
the partition boundary and merging the flanking exons. The pre-mRNA is
partitioned into a sequence of exon blocks $\{E_1, E_2, \ldots, E_N\}$
separated by intron partitions $\{I_1, I_2, \ldots, I_{N-1}\}$.
\end{definition}

\begin{theorem}[Splicing as Categorical Partition]
\label{thm:splicing_partition}
The probability of including exon $i$ in the mature mRNA is
\begin{equation}
  P(\text{include}\,E_i) = \frac{b^{-\Depth_{\mathrm{include},i}}}{\displaystyle\sum_{j}
  b^{-\Depth_{\mathrm{include},j}}},
  \label{eq:splicing_probability}
\end{equation}
where $\Depth_{\mathrm{include},i}$ is the partition depth of the
``include exon $i$'' path through the spliceosomal decision landscape, which
depends on:
\begin{enumerate}[label=(\alph*)]
  \item The intrinsic strength of the splice sites flanking $E_i$ (sequence
        information content, binding energy for U1/U2 snRNPs).
  \item Exon definition signals (ESEs, ESSs, ISEs, ISSs) that present or
        conceal the splice sites.
  \item The electrostatic context of the pre-mRNA, which modulates the
        binding affinity of SR proteins and hnRNPs to exonic and intronic
        regulatory elements.
\end{enumerate}
\end{theorem}

\begin{proof}
The spliceosomal reaction at each alternative splice site is a binding
competition: splice site $5'_{i}$ recruits U1 snRNP with rate
$k_{i}^{(5')} \propto \exp(-\Depth_{i}^{(5')}/\log_b e)$, and splice site
$3'_{i}$ recruits U2AF with rate
$k_{i}^{(3')} \propto \exp(-\Depth_{i}^{(3')}/\log_b e)$.
The overall depth for inclusion of exon $i$ is the sum over both splice site
decisions:
$\Depth_{\mathrm{include},i} = \Depth_i^{(5')} + \Depth_i^{(3')}$.
The partition selection rule from Part~I
(Chapter~\ref{chap:bounded_phase_space}) gives the inclusion probability as
the Boltzmann weight over all competing paths, which is
Eq.~\eqref{eq:splicing_probability}. $\blacksquare$
\end{proof}

%------------------------------------------------------------------------------
\subsection{Metabolic Coupling to Splicing}
\label{subsec:splicing_metabolic}
%------------------------------------------------------------------------------

Equation~\eqref{eq:splicing_probability} contains $\Depth_{\mathrm{include},i}$,
which depends on the electrostatic context of the pre-mRNA. This is the entry
point for metabolic regulation of splicing.

\begin{proposition}[Electrostatic Modulation of Splice Site Depth]
\label{prop:splice_depth_modulation}
The electrostatic component of $\Depth_{\mathrm{include},i}$ is
\begin{equation}
  \Depth_{\mathrm{include},i}^{(\mathrm{elec})} =
  -\frac{E_{\mathrm{bind}}(\lambda_D, \{q_{\mathrm{SR}}\})}{\kB T \ln b},
  \label{eq:splice_elec_depth}
\end{equation}
where $E_{\mathrm{bind}}$ is the binding energy of SR proteins to the exonic
splicing enhancer (ESE) of exon $i$, which depends on:
\begin{itemize}
  \item $\lambda_D$: the Debye length sets the range of electrostatic attraction
        between the positively charged RS domains of SR proteins and the
        negatively charged RNA backbone.
  \item $\{q_{\mathrm{SR}}\}$: the net charge of the SR protein RS domain,
        which is modified by phosphorylation (each phosphorylation adds $-2e$).
\end{itemize}
\end{proposition}

The key insight is that SR protein phosphorylation state is metabolically
regulated: SRPK1/2 (SR protein kinases) are activated by growth factor
signalling and nutrient availability, while PP1/PP2A phosphatases are
inhibited by high [ATP]. Under high-energy conditions, SR proteins are
hyperphosphorylated; under energy stress, they are hypophosphorylated.

\begin{corollary}[Metabolic Splicing Switch]
\label{cor:metabolic_splicing}
The splicing ratio of an alternatively spliced gene responds to metabolic state
through SR protein phosphorylation. Increasing metabolic activity
(higher ATP, higher SRPK activity) increases $|q_{\mathrm{SR}}|$
(more negative charge), which deepens $\Depth_{\mathrm{include},i}$ for
ESE-dependent exons (weaker binding of hyperphosphorylated SR proteins to
ESEs), thereby reducing the inclusion probability.
The converse holds under energy stress.
\end{corollary}

This mechanism explains a key empirical observation: perturbations of cellular
energy status (glucose withdrawal, hypoxia, oligomycin) alter splicing patterns
genome-wide within 5--15 minutes --- far faster than transcriptional
reprogramming could account for. The partition depth formula predicts these
changes because $\lambda_D$ and SR protein phosphorylation state both respond
to metabolic oscillations on sub-minute timescales.

\begin{keyresult}
The experimentally measured PKM1/PKM2 alternative splicing ratio oscillates
with 30\% amplitude at metabolic timescales. Partition theory accounts for this
oscillation through metabolic modulation of $\Depth_{\mathrm{include}}$ for
exons 9 and 10 of PKM via hnRNPA1/A2 and PTB phosphorylation state.
\end{keyresult}

%------------------------------------------------------------------------------
\section{Histone Modifications as Partition Depth Markers}
\label{sec:histone_modifications}
%------------------------------------------------------------------------------

The post-translational modification of histone tail residues --- acetylation,
methylation, phosphorylation, ubiquitylation, and others --- has long been
described as a ``histone code'' that is read by effector proteins to determine
chromatin state. Partition theory provides the first \emph{mechanistic}
interpretation of this code.

\begin{theorem}[Histone Modifications as Partition Depth Increments]
\label{thm:histone_code}
Each histone modification changes the partition depth $\Depth$ of the
associated chromatin region by a specific increment $\Delta\Depth$:
\begin{align}
  \text{H3K4me3 (active TSS):} &\quad \Delta\Depth = -\Depth_{\mathrm{rec}},
    \quad \Depth_{\mathrm{rec}} > 0,
    \label{eq:H3K4me3}\\
  \text{H3K36me3 (active gene body):} &\quad \Delta\Depth = -\Depth_{\mathrm{body}},
    \label{eq:H3K36me3}\\
  \text{H3K27me3 (Polycomb repression):} &\quad \Delta\Depth = +\Depth_{\mathrm{PRC2}},
    \label{eq:H3K27me3}\\
  \text{H3K9me3 (constitutive heterochromatin):} &\quad \Delta\Depth = +\Depth_{\mathrm{HP1}},
    \label{eq:H3K9me3}
\end{align}
where active marks carry negative increments (decreasing depth, more
accessible) and repressive marks carry positive increments (increasing depth,
less accessible). The modifications are not the primary cause of chromatin
state; they are \emph{records of the partition depth at the time of writing}
--- stable annotations that allow the depth state to persist through
cell division and perturbation.
\end{theorem}

\begin{proof}
H3K4me3 is deposited by SETD1A/B and MLL complexes, which are recruited by
the CTD of RNA Pol~II in its Ser5-phosphorylated (initiating) form. The
modification appears co-transcriptionally and correlates with low nucleosome
occupancy --- \ie, low $\Depth$. The mark is a consequence of the low-$\Depth$
state, not its cause. Similarly, H3K27me3 is deposited by PRC2, which is
recruited to compact, high-$\Depth$ chromatin and uses allosteric activation
by H3K27me3 itself to spread. The modification amplifies and records a
pre-existing high-$\Depth$ state. Formally, the modification acts as a
positive feedback term in the $\Depth$ evolution equation:
$\dot\Depth_{\mathrm{mod}} = \alpha\,\Depth + \beta\,[\text{modifier}]$,
driving $\Depth$ to either 0 (active, low-$\Depth$ attractor) or
$\Depth_{\max}$ (silent, high-$\Depth$ attractor). $\blacksquare$
\end{proof}

\begin{definition}[The Partition Code]
\label{def:partition_code}
The histone code is the partition code: the set of chromatin modifications
that record the partition depth state of each genomic locus in a form that is
heritable through replication (via reader--writer coupling) and readable by
effector proteins (via modification-specific binding domains). The partition
code is not combinatorial in an arbitrary sense; it is a scalar record of
$\Depth$ at each position.
\end{definition}

%------------------------------------------------------------------------------
\subsection{DNA Methylation as an Additional Depth Increment}
\label{subsec:dna_methylation}
%------------------------------------------------------------------------------

DNA methylation at CpG dinucleotides (5-methylcytosine, 5mC) adds an
additional partition depth increment that is independent of the histone
modification state:

\begin{proposition}[5mC Depth Increment]
\label{prop:5mc_depth}
Methylation of a CpG site increases the partition depth of the overlapping
chromatin region by $\Delta\Depth_{\mathrm{5mC}}$, arising from two mechanisms:
\begin{enumerate}[label=(\roman*)]
  \item Direct steric obstruction of CpG-reading transcription factors (TFBSs
        containing CpG are occluded by the methyl group).
  \item Recruitment of methyl-CpG binding proteins (MeCP2, MBD1--4, KAISO),
        which are scaffolds for HDAC-containing repressor complexes, further
        increasing $\Depth$ through histone deacetylation.
\end{enumerate}
The combined depth increment is:
\begin{equation}
  \Delta\Depth_{\mathrm{5mC}} = \Delta\Depth_{\mathrm{steric}}
  + \Delta\Depth_{\mathrm{HDAC}},
  \label{eq:methylation_depth}
\end{equation}
making 5mC a doubly reinforced partition depth increment.
\end{proposition}

%------------------------------------------------------------------------------
\section{The Chromatin State Machine}
\label{sec:chromatin_state_machine}
%------------------------------------------------------------------------------

Combining nucleosome breathing (Section~\ref{sec:breathing_overview}),
partition gate action (Section~\ref{sec:nucleosome_gate}), and histone
modification depth increments (Section~\ref{sec:histone_modifications})
yields a complete picture of chromatin as a state machine.

\begin{definition}[Chromatin State Machine]
\label{def:chromatin_state_machine}
A chromatin locus at position $x$ is described by the state vector
\begin{equation}
  \mathbf{C}(x,t) = \bigl(\Depth(x,t),\,\theta(x,t),\,\mathcal{M}(x)\bigr),
  \label{eq:chromatin_state}
\end{equation}
where $\Depth(x,t)$ is the instantaneous partition depth, $\theta(x,t)$ is the
nucleosome breathing angle, and $\mathcal{M}(x)$ is the vector of modification
marks (a discrete coordinate). The state evolves according to:
\begin{align}
  \dot\Depth(x,t) &= -k_{\mathrm{rem}}\Depth + k_{\mathrm{mod}}\mathcal{M}
    + k_\lambda(\lambda_D(t) - \bar\lambda_D), \label{eq:depth_evolution}\\
  \dot\theta(x,t) &= -k_\theta(\theta - \theta_{\mathrm{eq}}(\lambda_D(t))),
    \label{eq:theta_evolution}
\end{align}
where $\theta_{\mathrm{eq}}(\lambda_D)$ is the equilibrium breathing angle at
Debye length $\lambda_D$ (from Theorem~\ref{thm:nucleosome_breathing}), and
$k_\lambda$ captures the rate of depth change induced by Debye length
deviations from the mean.
\end{definition}

The chromatin state machine is driven by the metabolic clock through
$\lambda_D(t)$ in Eqs.~\eqref{eq:depth_evolution}--\eqref{eq:theta_evolution}.
Its stationary states are the chromatin states observed by ChIP-seq,
ATAC-seq, and Hi-C: active promoters ($\Depth \approx 0$, $\theta \approx 0.8$),
Polycomb-repressed genes ($\Depth \approx \Depth_{\mathrm{PRC2}}$,
$\theta \approx 0.1$), and constitutive heterochromatin
($\Depth \approx \Depth_{\max}$, $\theta \approx 0$).

%------------------------------------------------------------------------------
\section{Experimental Predictions}
\label{sec:ch19_predictions}
%------------------------------------------------------------------------------

\begin{proposition}[Nucleosome Breathing Amplitude Scales with Debye Length Oscillation]
\label{prop:breathing_scaling}
Perturbations that increase $\delta\lambda_D/\bar\lambda_D$ (e.g., reducing
extracellular \ce{K^+} to reduce intracellular ionic strength) should increase
nucleosome breathing amplitude proportionally (Eq.~\eqref{eq:breathing_amplitude}).
FRET-based single-nucleosome measurements in live cells should detect this
scaling relationship.
\end{proposition}

\begin{proposition}[SR Protein Phosphorylation Predicts Splicing Ratios]
\label{prop:SR_splicing}
Across cell types with different metabolic rates, the phosphorylation ratio of
SR proteins (measured by phospho-proteomics) should predict the PKM1/PKM2
splicing ratio more accurately than SR protein total abundance, because
phosphorylation sets $q_{\mathrm{SR}}$ and therefore $\Depth_{\mathrm{include}}$
in Eq.~\eqref{eq:splicing_probability}.
\end{proposition}

\begin{proposition}[Histone Modification Writing is Delayed Relative to Transcription]
\label{prop:histone_writing_delay}
Because histone modifications record the depth state rather than creating it,
they should appear with a measurable delay after changes in transcription.
ChIP-seq time courses following acute transcriptional induction should show
H3K4me3 acquisition peaking 15--30\,min after RNA Pol~II engagement, not
simultaneously.
\end{proposition}

%------------------------------------------------------------------------------
\begin{chapterbox}
\textbf{Chapter Summary.}
\begin{itemize}
  \item Nucleosome breathing is driven by Debye length oscillations
        (Theorem~\ref{thm:nucleosome_breathing}): the 50\% amplitude
        oscillation at metabolic timescales is the quantitative signature of
        electrostatic coupling.
  \item A nucleosome is a partition gate (Theorem~\ref{thm:nucleosome_gate}):
        closed state sets $\Depth_{\mathrm{access}} \to \infty$; open state
        reduces $\Depth_{\mathrm{access}}$ to allow factor binding.
        Chromatin remodellers are partition topology modifiers.
  \item Alternative splicing is categorical partition selection
        (Theorem~\ref{thm:splicing_partition}):
        $P(\text{include}\,E_i) = b^{-\Depth_{\mathrm{include},i}}/
        \sum_j b^{-\Depth_{\mathrm{include},j}}$.
        Metabolic perturbations change splicing ratios within minutes through
        SR protein phosphorylation state (Corollary~\ref{cor:metabolic_splicing}).
  \item The histone code is the partition code
        (Theorem~\ref{thm:histone_code}): active marks are negative depth
        increments; repressive marks are positive depth increments.
        5mC adds an additional doubly reinforced depth increment.
  \item The chromatin state machine
        (Definition~\ref{def:chromatin_state_machine}) is driven by the
        metabolic clock through the Debye length oscillation.
\end{itemize}
\end{chapterbox}
